\documentclass[11pt, a4paper]{article}

\usepackage[a4paper, margin=2.5cm]{geometry}
\usepackage[utf8]{inputenc}
\usepackage[T1]{fontenc}
% lmodern not available in this environment
% microtype disabled
\usepackage{parskip}
\usepackage{setspace}
\usepackage{titlesec}
\usepackage{booktabs}
\usepackage{array}
\usepackage{tabularx}
\usepackage{xcolor}
\usepackage{amsmath}
\usepackage{amssymb}
\usepackage{enumitem}
\usepackage{hyperref}
\usepackage{fancyhdr}
\usepackage{mdframed}
\usepackage{xurl}

\definecolor{darkblue}{HTML}{1a1a2e}
\definecolor{midblue}{HTML}{2e2e5e}
\definecolor{lightbg}{HTML}{f0f0fa}
\definecolor{framebg}{HTML}{eef0ff}
\definecolor{framebd}{HTML}{aaaacc}
\definecolor{mathbg}{HTML}{f5f5ff}

\hypersetup{
    colorlinks=true,
    linkcolor=midblue,
    urlcolor=midblue,
    pdftitle={Preprocessing Pipeline --- Technical Explanation},
}

\titleformat{\section}
  {\large\bfseries\color{darkblue}}
  {\thesection.}{0.6em}{}
  [\vspace{-0.3em}\textcolor{framebd}{\rule{\linewidth}{0.4pt}}]

\titleformat{\subsection}
  {\normalsize\bfseries\color{midblue}}
  {\thesubsection}{0.6em}{}

\titlespacing{\section}{0pt}{18pt}{6pt}
\titlespacing{\subsection}{0pt}{10pt}{4pt}

\pagestyle{fancy}
\fancyhf{}
\renewcommand{\headrulewidth}{0.4pt}
\fancyhead[L]{\small\textcolor{darkblue}{\textbf{Preprocessing Pipeline}}}
\fancyhead[R]{\small\textcolor{darkblue}{201940185}}
\fancyfoot[C]{\small\thepage}

\newmdenv[
  backgroundcolor=mathbg,
  linecolor=framebd,
  linewidth=0.6pt,
  innerleftmargin=12pt,
  innerrightmargin=12pt,
  innertopmargin=8pt,
  innerbottommargin=8pt,
  skipabove=6pt,
  skipbelow=6pt,
]{mathbox}

\newmdenv[
  backgroundcolor=framebg,
  linecolor=framebd,
  linewidth=0.6pt,
  innerleftmargin=12pt,
  innerrightmargin=12pt,
  innertopmargin=8pt,
  innerbottommargin=8pt,
  skipabove=6pt,
  skipbelow=6pt,
]{insightbox}

\setlist[itemize]{leftmargin=1.4em, itemsep=3pt, topsep=4pt}

\begin{document}

\begin{center}
  {\LARGE\bfseries\color{darkblue} Preprocessing Pipeline}\\[4pt]
  {\normalsize CA Assignment 2 $\cdot$ Clustering Algorithms}\\[2pt]
  {\normalsize Embedding model:
    \texttt{sentence-transformers/all-mpnet-base-v2}}
\end{center}

\vspace{4pt}
\textcolor{framebd}{\rule{\linewidth}{0.8pt}}
\vspace{6pt}

\section*{Overview}

The preprocessing pipeline transforms raw natural language sentences into a
compact, geometrically well-structured numerical representation suitable for
Agglomerative Hierarchical Clustering (AHC).  The pipeline consists of five
sequential stages:

\begin{enumerate}
    \item Data loading and text cleaning
    \item Sentence embedding using a pre-trained transformer model
    \item L2 normalisation of the raw embedding vectors
    \item Adaptive dimensionality reduction via Principal Component Analysis (PCA)
    \item Non-linear manifold projection via UMAP
\end{enumerate}

All 1{,}760 rows are processed, including duplicate sentences.  This is a
deliberate design decision explained in full below.

\section{Data Loading and Text Cleaning}

The dataset is stored in a tab-separated file with two columns: an integer
identifier (\texttt{ID}) and a raw sentence string (\texttt{Sentence}).  The
file is read line by line using a custom parser that identifies record
boundaries by detecting lines whose first field is a positive integer followed
by a tab character.  Any line that does not match this pattern is treated as a
continuation of the previous sentence.  This approach correctly handles
sentences that span multiple physical lines, ensuring all 1{,}760 instances
are loaded.

\subsection{Cleaning Strategy}

A deliberate \textbf{light-touch cleaning strategy} is adopted:

\begin{itemize}
  \item Removal of escaped quotation marks (\verb|\"| and \verb|\'|).
  \item Collapsing of repeated whitespace characters into a single space.
  \item Collapsing of actual newline characters into a single space.
  \item Stripping of leading and trailing whitespace.
  \item Replacement of empty strings with the placeholder \texttt{`empty'}.
\end{itemize}

Heavy-handed normalisation (lowercasing, punctuation removal, stopword
removal) is intentionally excluded.  \texttt{all-mpnet-base-v2} is trained on
natural, cased, punctuated English text.  Its self-attention mechanism is
sensitive to capitalisation (e.g., character prefixes such as
\textsc{Coriolanus:} carry strong identity signals), punctuation, and function
words.  Removing these would push sentences out of the model's training
distribution, degrading embedding quality.

\subsection{Duplicate Retention}

The dataset contains 64 duplicate rows: 18 unique Notre Dame / Wikipedia
sentences each repeated 4--5 times, giving 80 rows total.  These duplicates
are \textbf{intentionally retained} rather than removed.

The reason is geometric: duplicate sentences produce identical embedding
vectors, which map to the same point in UMAP space.  With \texttt{n\_neighbors=10},
UMAP constructs a $k$-nearest-neighbour graph in which each Notre Dame
sentence finds most or all of its 10 nearest neighbours within the Notre Dame
group itself.  This creates a densely connected sub-graph that UMAP projects
into a tight, isolated manifold region, which AHC then detects as a distinct
cluster.

Without the duplicate rows, the 18 unique Notre Dame sentences represent only
1.1\% of the dataset.  At that density, each sentence's 10-nearest-neighbour
graph would be dominated by non-Notre-Dame neighbours, the sub-graph would not
be cohesive, and the Notre Dame region would dissolve into the surrounding
modern English manifold.

\begin{insightbox}
\textbf{Influence on clustering:} Retaining duplicates is the mechanism that
enables the detection of the Notre Dame sub-domain as a geometrically
distinct cluster.  Removing them would merge all Notre Dame content into the
broader modern English cluster, reducing the clustering to $K=2$.
\end{insightbox}

\section{Sentence Embedding with all-mpnet-base-v2}

Sentences are encoded into dense numerical vectors using the pre-trained model
\textbf{\path{sentence-transformers/all-mpnet-base-v2}}, a 12-layer,
\textbf{768-dimensional} sentence transformer based on the MPNet architecture
(Song et al., 2020).  The model was fine-tuned on over one billion sentence
pairs using a Multiple Negative Ranking (MNR) loss, which optimises the
embedding space so that semantically similar sentences map to geometrically
proximate points on the unit hypersphere in $\mathbb{R}^{768}$.

\subsection{Why This Model?}

\begin{itemize}
  \item \textbf{Highest-quality standard sentence transformer.}
        \texttt{all-mpnet-base-v2} achieves the highest overall score among
        the standard \texttt{all-*} family on the Sentence Embeddings
        Benchmark and on MTEB clustering tasks.  Petukhova et al.\ (2024)
        confirm BERT-family transformers are the best open-source option for
        text clustering.

  \item \textbf{768-dimensional output.}  With 110 million parameters and
        twice the embedding dimensions of MiniLM-L12-v2, the model encodes
        substantially finer-grained semantic distinctions --- essential for
        separating four linguistically diverse sub-corpora.

  \item \textbf{Cosine-similarity training objective.}  The MNR loss trains
        cosine similarity to be proportional to semantic similarity.  This
        makes the embedding space intrinsically well-suited to cosine-based
        distance measures, which are used in all subsequent steps.

  \item \textbf{No cloud API dependency.}  The model runs fully locally via
        Hugging Face's model hub.
\end{itemize}

\subsection{Encoding Configuration}

Sentences are encoded in batches of 64 with
\texttt{normalize\_embeddings=False}.  Normalisation is applied explicitly in
the next step.  The output is a \texttt{float32} matrix of shape
$(1760,\,768)$.

\begin{insightbox}
\textbf{Influence on clustering:} Sentences from semantically distinct
registers --- Shakespearean dialogue, contemporary reviews, and encyclopaedic
prose --- are embedded in clearly separated regions of $\mathbb{R}^{768}$,
providing the raw geometric signal from which clusters are recovered.
\end{insightbox}

\section{L2 Normalisation}

Each raw embedding vector $\mathbf{x} \in \mathbb{R}^{768}$ is L2-normalised:

\begin{mathbox}
\[
  \hat{\mathbf{x}} \;=\; \frac{\mathbf{x}}{\|\mathbf{x}\|_2}
  \qquad \text{such that} \qquad
  \|\hat{\mathbf{x}}\|_2 = 1
\]
\end{mathbox}

\subsection{Why Normalise Before PCA?}

Without normalisation, embeddings vary in both direction (semantic content) and
magnitude (loosely correlated with sentence length).  PCA applied to
un-normalised vectors decomposes a mixture of directional and magnitude
variance.  L2 normalisation ensures PCA decomposes purely
\textbf{angular (semantic) variance}.

A critical geometric consequence is that, for unit vectors
$\hat{\mathbf{a}}, \hat{\mathbf{b}}$:

\begin{mathbox}
\[
  d_{\mathrm{Euclidean}}(\hat{\mathbf{a}},\,\hat{\mathbf{b}})^{2}
  \;=\;
  2\bigl(1 - \cos(\hat{\mathbf{a}},\, \hat{\mathbf{b}})\bigr)
\]
\end{mathbox}

Euclidean distance is a \textbf{monotone function of cosine distance} on the
unit sphere.  This ensures that UMAP's cosine neighbourhood graph and
Ward-linkage AHC (Euclidean) are both consistent with the MNR training
objective of \texttt{all-mpnet-base-v2}.

\begin{insightbox}
\textbf{Influence on clustering:} L2 normalisation aligns the distance metric
used in every subsequent step (cosine for UMAP neighbourhood graph, Euclidean
for PCA projection and AHC) with the metric optimised during model training,
ensuring geometric coherence throughout the pipeline.
\end{insightbox}

\section{Adaptive Principal Component Analysis}

PCA is applied to the L2-normalised matrix $\hat{X} \in \mathbb{R}^{1760
\times 768}$.  Rather than fixing a component count, an
\textbf{adaptive} strategy is used: the minimum number of components $n^*$
whose cumulative explained variance reaches threshold $\tau = 0.80$ is
selected:

\begin{mathbox}
\[
  n^{*} \;=\; \min\!\left\{n \;\middle|\;
    \frac{\sum_{k=1}^{n} \lambda_k}{\sum_{k=1}^{p} \lambda_k}
    \;\geq\; \tau \right\},
  \qquad \tau = 0.80,\quad p = 768
\]
\end{mathbox}

For this dataset, $n^{*} = 121$ components are retained, preserving
\textbf{80.02\%} of total variance.

\subsection{Why Apply PCA?}

\begin{itemize}
  \item \textbf{Curse of dimensionality.}  In $\mathbb{R}^{768}$, pairwise
        Euclidean distances become nearly uniform (Beyer et al., 1999),
        degrading distance-based algorithms.  Reducing to 121 dimensions
        substantially mitigates this.

  \item \textbf{Denoising.}  Components 122 through 768 capture
        sentence-level idiosyncrasies rather than shared semantic structure.
        Discarding them acts as a low-rank denoising filter.

  \item \textbf{UMAP pre-conditioning.}  UMAP's $k$-nearest-neighbour graph
        construction is $O(N \log N)$ in the embedding dimension.  Reducing
        from 768 to 121 dimensions substantially reduces this cost while
        preserving the semantic structure that UMAP needs to resolve the
        manifold.
\end{itemize}

\subsection{Why Not Re-Normalise After PCA?}

PCA is variance-preserving: PC$_1$ carries the largest eigenvalue, PC$_2$ the
second largest, and so on.  Euclidean distance in PCA space naturally weights
the earlier, higher-variance components:

\begin{mathbox}
\[
  d^{2}(\mathbf{a},\,\mathbf{b})
  \;=\;
  \sum_{k=1}^{121}(a_k - b_k)^{2}
\]
\end{mathbox}

This variance weighting provides useful inductive bias for UMAP: the most
semantically dominant axes (e.g.\ the contrast between archaic and modern
English) automatically receive greater influence in the neighbourhood graph.
Re-normalising after PCA would collapse all 121 dimensions to equal
contribution, erasing this structure.

\begin{insightbox}
\textbf{Influence on clustering:} Adaptive PCA retains exactly the components
needed to represent 80\% of the semantic variance in the 768-dimensional
embedding space, providing UMAP with a clean, noise-reduced,
variance-weighted input.
\end{insightbox}

\section{UMAP: Non-Linear Manifold Projection}

After PCA denoising, Uniform Manifold Approximation and Projection (UMAP;
McInnes, Healy \& Melville, 2018) projects the 121-dimensional PCA
coordinates into a 5-dimensional Euclidean space used directly for clustering.

\subsection{Why UMAP After PCA?}

PCA is linear: it decomposes global variance along orthogonal axes and cannot
represent non-linear cluster structure.  Sentence embeddings lie on a curved
manifold in $\mathbb{R}^{768}$; sub-domains that are genuinely distinct may
not be linearly separable.  UMAP constructs a fuzzy topological representation
of the data manifold using a $k$-nearest-neighbour graph, then optimises a
low-dimensional layout that preserves both local neighbourhood structure and
global topology simultaneously.

\subsection{Parameter Choices}

\begin{itemize}
  \item \textbf{\texttt{metric='cosine'}.}  UMAP builds its $k$-nearest-neighbour
        graph using cosine distance in the PCA-denoised space.  Cosine
        distance is the semantically correct measure for L2-normalised text
        embeddings, consistent with the MNR training objective.  This is the
        BERTopic standard (Grootendorst, 2022).

  \item \textbf{\texttt{n\_components=5}.}  Standard output dimensionality for
        clustering tasks (BERTopic standard).  Sufficient to preserve
        multi-cluster manifold topology; low enough for Euclidean distances to
        remain well-conditioned for AHC.

  \item \textbf{\texttt{min\_dist=0.0}.}  The UMAP documentation states
        explicitly that low values are recommended for clustering tasks.
        \texttt{min\_dist=0.0} allows points within the same neighbourhood to
        pack as tightly as the manifold permits, producing compact, isolated
        cluster regions.  This is the BERTopic standard for clustering.

  \item \textbf{\texttt{n\_neighbors=10}.}  Reduced from the standard default
        of 15 to increase sensitivity to small, dense local structures.  With
        \texttt{n\_neighbors=10}, each of the 80 Notre Dame rows finds most of
        its 10 nearest neighbours within the Notre Dame group itself, creating
        a strongly connected sub-graph that UMAP projects into a tight,
        isolated manifold region.  The BERTopic documentation explicitly states
        that decreasing \texttt{n\_neighbors} ``creates more local structure''
        and is appropriate when detecting small, cohesive sub-domains.

  \item \textbf{\texttt{random\_state=42}.}  Fixes UMAP's stochastic gradient
        descent phase for bit-for-bit reproducibility.
\end{itemize}

\begin{insightbox}
\textbf{Influence on clustering:} UMAP with \texttt{n\_neighbors=10} is the
key step that enables the detection of the Notre Dame sub-domain as a
geometrically distinct cluster.  The reduced neighbourhood size makes the
algorithm more sensitive to the dense, localised cluster formed by the
repeated Notre Dame sentences, which UMAP projects into a compact isolated
region that Ward-linkage AHC can cleanly separate from the surrounding
modern English manifold.
\end{insightbox}

\section{Pipeline Summary}

\begin{table}[h!]
\centering
\renewcommand{\arraystretch}{1.35}
\begin{tabularx}{\textwidth}{c l l l X}
  \toprule
  \textbf{Stage} & \textbf{Operation} & \textbf{Input}
                 & \textbf{Output} & \textbf{Key Purpose} \\
  \midrule
  1 & Data loading \& cleaning
    & Raw text file      & 1{,}760 sentences (dups retained)
    & Quality \& row alignment \\
  2 & all-mpnet-base-v2 encoding
    & 1{,}760 sentences  & $(1760,\,768)$
    & Semantic representation \\
  3 & L2 normalisation
    & $(1760,\,768)$     & $(1760,\,768)$ unit vectors
    & Cosine-consistent distances \\
  4 & Adaptive PCA (121 components)
    & $(1760,\,768)$     & $(1760,\,121)$ variance-weighted
    & Denoising; UMAP pre-conditioning \\
  5 & UMAP ($n_{\text{nbrs}}=10$, $\texttt{min\_dist}=0$)
    & $(1760,\,121)$     & $(1760,\,5)$ manifold-projected
    & Local cluster separation \\
  \bottomrule
\end{tabularx}
\caption{Summary of the five-stage preprocessing pipeline.}
\end{table}

\section{Holistic Discussion: Influence on Clustering Performance}

The preprocessing choices form a coherent, theoretically motivated pipeline:

\begin{itemize}
  \item \textbf{Cleaning $\to$ Embedding quality.}  Preserving natural text
        keeps sentences within the model's training distribution.  Archaic
        Shakespeare vocabulary, contemporary review prose, and formal
        encyclopaedic text produce embeddings in clearly separated regions
        of $\mathbb{R}^{768}$.

  \item \textbf{Duplicate retention $\to$ Notre Dame cluster density.}
        The 80 Notre Dame rows (18 unique $\times$ 4--5 copies) create a
        dense local neighbourhood in UMAP space.  This is the necessary
        precondition for AHC to detect the Notre Dame sub-domain as a
        distinct cluster.

  \item \textbf{L2 normalisation $\to$ Metric consistency.}  All steps from
        PCA onwards operate in cosine-equivalent Euclidean space, consistent
        with the MNR training objective.

  \item \textbf{Adaptive PCA $\to$ UMAP pre-conditioning.}  Noise suppression
        and dimensionality reduction produce a clean, variance-weighted input
        that allows UMAP to resolve the non-linear manifold structure
        efficiently.

  \item \textbf{UMAP ($n_{\text{nbrs}}=10$) $\to$ Local cluster isolation.}
        The reduced neighbourhood size amplifies local density differences,
        enabling the detection of the small but cohesive Notre Dame cluster
        that would be invisible to a linear projection or to UMAP with larger
        neighbourhood sizes.
\end{itemize}

\vspace{12pt}
\noindent\textcolor{midblue}{\small\textbf{References:}\\[3pt]
Beyer, K., Goldstein, J., Ramakrishnan, R., \& Shaft, U.\ (1999).\ When is
``nearest neighbor'' meaningful? \textit{ICDT}, 217--235.\\[3pt]
Grootendorst, M.\ (2022).\ BERTopic: Neural topic modelling with a class-based
TF-IDF procedure. \textit{arXiv:2203.05794}.\\[3pt]
McInnes, L., Healy, J., \& Melville, J.\ (2018).\ UMAP: Uniform manifold
approximation and projection for dimension reduction.
\textit{arXiv:1802.03426}.\\[3pt]
Petukhova, A., Matos-Carvalho, J.\,P., \& Fachada, N.\ (2024).\ Text clustering
with LLM embeddings. \textit{arXiv:2403.15112}.\\[3pt]
Reimers, N., \& Gurevych, I.\ (2019).\ Sentence-BERT: Sentence embeddings using
Siamese BERT-networks. \textit{Proceedings of EMNLP 2019}, 3982--3992.\\[3pt]
Song, K., Tan, X., Qin, T., Lu, J., \& Liu, T.-Y.\ (2020).\ MPNet: Masked and
permuted pre-training for language understanding.
\textit{NeurIPS}, 33, 16857--16867.}

\end{document}